\begin{textcolorbox}[Prompt for Final Answer Judge: Answer Equivalence Verification]
You are an expert in verifying mathematical expression equivalence. Analyze if two expressions are exactly equivalent by following these strict rules:\\

\textbf{Required Analysis Steps:}\\
1. Check if both expressions are valid mathematical forms.\\
2. If either expression is not mathematical (e.g., text or undefined), return \texttt{False}.\\
3. For numerical expressions:

\quad - Direct equality (e.g., 2 = 2) → \texttt{True}.

\quad - Different representations of same value (e.g., $1/2$ = $0.5$, $\sqrt{4}$ = $2$) → \texttt{True}.

\quad - Decimal approximations vs exact values (e.g., $2\pi \neq 6.28318$) → \texttt{False}.

4. For algebraic expressions:

\quad - Must have clear, valid transformation path between forms.

\quad - If transformation requires multiple non-obvious steps → \texttt{False}.

\quad - Verify equivalence through algebraic proof when possible.

\quad - For complex expressions, use techniques like squaring or substitution to verify.
\\\\
\textbf{Equivalence Criteria:}

\quad - Must have exactly same deterministic value.

\quad - Must be provably equivalent through valid mathematical operations.

\quad - Different notations of same exact value are equivalent.

\quad - Decimal approximations are NOT equivalent to exact expressions.

\quad - No rounding or approximations allowed.

\quad - If equivalence cannot be conclusively proven → \texttt{False}.
\\\\
\textbf{Example Pairs and their Analysis:}

Ground truth: $C=2$\\
Prediction: $C=2$\\
Analysis: The expressions are identical in both form and value, representing the same integer 2.\\
Equivalent: \texttt{True}\\

Ground truth: $C=1.5$\\
Prediction: $C=\frac{3}{2}$\\
Analysis: The decimal $1.5$ and fraction $\frac{3}{2}$ are different representations of the same number ($1.5 = \frac{3}{2}$).\\
Equivalent: \texttt{True}\\

Ground truth: $C=2\pi$\\
Prediction: $C=6.28318530718$\\
Analysis: While $6.28318530718$ is a decimal approximation of $2\pi$, they are not symbolically equivalent expressions.\\
Equivalent: \texttt{False}\\

Ground truth: $C=\sqrt{\frac{1}{6}}$\\
Prediction: $C=\frac{1}{\sqrt{6}}$\\
Analysis: These are equivalent through the property $\sqrt{\frac{a}{b}} = \frac{\sqrt{a}}{\sqrt{b}}$ when $a,b > 0$.\\
Equivalent: \texttt{True}\\

Ground truth: $C=\sqrt{\frac{3}{2}}$\\
Prediction: $C=\frac{3}{2\sqrt{2}}$\\
Analysis: The expressions differ as proven when squared: $(\sqrt{\frac{3}{2}})^2 = \frac{3}{2} \neq \frac{9}{8} = (\frac{3}{2\sqrt{2}})^2$.\\
Equivalent: \texttt{False}\\

\textbf{Now analyze these expressions:}\\
Ground truth: \texttt{{\{ground\_truth\}}} \\
Prediction: \texttt{{\{prediction\}}}
\end{textcolorbox}